% short/03_physics.tex: why the hand is detectable (about 1.25 pages with Figure 4). Full version:
% sections/03_physics.tex (eq. mirror kept; tab:mirror and the class-means figure left to the full report).
\section{Why the hand is detectable}\label{sec:physics}

\paragraph{The mirror.} Let $\vf$ be the specific force (the browser's acceleration including gravity), $\va$ the
acceleration and $\vomega$ the angular velocity, all in device coordinates. Mirroring the body across its sagittal
plane and relabelling the device's $x$ axis gives, with $\Smat = \operatorname{diag}(-1, 1, 1)$, for every grip and
every movement
\begin{equation}
  \vf' = \Smat\vf, \qquad \va' = \Smat\va, \qquad \vomega' = -\Smat\vomega .
  \label{eq:mirror}
\end{equation}

\paragraph{Chiral quantities.} Angular velocity is a pseudovector, so with $\ghat$ the unit vector pointing up, the
yaw rate about gravity, $\vomega\cdot\ghat$, changes sign under the mirror and is unchanged by any rotation of the
phone in the hand. Only such chiral (mirror-odd) quantities can carry the hand: a mirror-even quantity has the same
distribution for both hands whenever the two hands move as mirror images.

\paragraph{The mechanism.} Early in the approach the forearm swings toward the body's midline, so, seen from above,
the phone turns counter-clockwise in the right hand and clockwise in the left. Near the nose the wrist turns the
phone partly back, so the early turn carries the clearest signal.

\begin{figure}[htbp]
\centering
\genfig{core/S2_13_first_half_yaw}
\caption{The yaw about gravity over the first half of the approach: AUC within each person (left) and the pooled
distributions (right). Its pooled AUC is higher than that of the net yaw (no paired test), so the models get
time-resolved inputs.}
\label{fig:S2_13_first_half_yaw}
\end{figure}

\paragraph{The evidence.} The yaw about gravity over the first half of the approach separates the hands with a
pooled area under the ROC curve (AUC) of \R{describe/fhy_auc_pooled}, against \R{describe/net_yaw_auc_pooled} for
the net yaw over the whole approach, a difference not tested in a paired comparison
(Figure~\ref{fig:S2_13_first_half_yaw}). Its within-person AUC is above 0.5 for
\R{describe/fhy_persons_auc_above_half} of \R{describe/n_persons} persons, its sign alone, with no training,
classifies \R{describe/fhy_sign_rule_accuracy} of the trials, and it separates the hands with the same sign on both
platforms (AUC \R{describe/fhy_auc_android} on Android, \R{describe/fhy_auc_ios} on iOS).

\paragraph{Chiral against mirror-even features.} In the feature catalogue of the interpretable models, chiral
features reach a LOSO AUC of up to \R{m1/cat_odd_auc_max}, while every mirror-even feature stays between
\R{m1/cat_even_auc_min} and \R{m1/cat_even_auc_max} (full report, Appendix~\fullref{app:m1:catalogue}). The two
hands are close to, but not exactly, mirror images, so mirror augmentation (a mirrored trial counted as a trial of
the other hand) is a plausible, not an exact, way to share data between them.

\paragraph{Consequences.} The models therefore get time-resolved chiral inputs, a body frame whose chiral channels
do not depend on the grip, and mirror augmentation.
The return of the arm is a time-reversed approach: on the \R{describe/return_trials} trials with a clean return
window, the first-half yaw rule is right on \R{describe/return_rule_acc_approach} of the approaches but on only
\R{describe/return_rule_acc_return} of the returns, so a cue-free detector that scores the return answers the wrong
hand.
